
\renewcommand{\Folder}{2025_MapaMuniciosMexico}
\renewcommand{\EntradaBibtex}{2025_MapaMuniciosMexico}

\begin{frame}{\citetitle{\EntradaBibtex} \footnotemark[1] (1)}

\begin{columns}
	\begin{column}{0.98\textwidth}
		\begin{itemize}
			\item Se tomaron como referencia archivos SHP obtenidos de CONABIO, process them using
			\item Se procesaron con Python y se los datos se integraton en una aplicación de Android utilizando JetJetpack Compose y Google Maps.
            
		\end{itemize}
	\end{column}
\end{columns}
%\end{block} 
\footnotetext[1]{\fullcite{\EntradaBibtex}}
\end{frame}


\begin{frame}{\citetitle{\EntradaBibtex}  (2)}

\begin{columns}
	\begin{column}{0.48\textwidth}
		\begin{itemize}
			\item Se tiene información de más de 2,400 municipios, por lo que su carga en el mapa puede ralentizar la app.
			\item Se utilizaron despachadores para cargar los archivos de datos individuales solo cuando el usuario da clic en el mapa.
		\end{itemize}
	\end{column}
	\begin{column}{0.2\textwidth}
		\begin{center}
			\includegraphics[width=0.98\textwidth]{\Folder/figs/captura1}
		\end{center}
	\end{column}
	
	\begin{column}{0.2\textwidth}
		 \begin{center}
		     \includegraphics[width=0.98\textwidth]{\Folder/figs/captura2}
		 \end{center}
	\end{column}
	\begin{column}{0.2\textwidth}
		 \begin{center}
		     \includegraphics[width=0.98\textwidth]{\Folder/figs/captura3}
		 \end{center}
	\end{column}
\end{columns}
%\end{block} 
%\footnotetext[1]{\fullcite{\EntradaBibtex}}
\end{frame}

